% ==============================================================================
% 03_methodik.tex – Kapitel 3
% ==============================================================================
\chapter{Methodik}
\label{ch:methodik}

Dieses Kapitel gibt einen Überblick über das Vorgehen: Daten
(\cref{sec:meth_daten,sec:meth_aufnahme}), Vorverarbeitung
(\cref{sec:meth_vorverarbeitung}), Modell und Training
(\cref{sec:meth_modell}), Auswertung (\cref{sec:meth_evaluation}) und
Implementierung (\cref{sec:meth_implementierung}). Die genauen Einstellungen
der Experimente stehen in \cref{sec:setup}.

\section{Datensätze}
\label{sec:meth_daten}
Die Arbeit nutzt zwei Datenquellen:
\begin{itemize}
  \item \textbf{GazeCapture}~\cite{krafka2016itracker} für die Basismodelle.
        Der vollständige Datensatz hat \num{2445504} Bilder von \num{1474}
        Personen. Verwendet wird eine öffentlich verfügbare Teilmenge von
        Kaggle~\cite{kaggle_gazecapture}. Das Skript \code{main.py} lädt sie
        herunter und liest für jede Person die Dateien \code{frames.json}
        (Bildnamen) und \code{dotInfo.json} (Position des Punktes auf dem
        Bildschirm, \code{XPts}/\code{YPts}). Personen mit fehlenden Dateien
        oder ungleich vielen Bildern und Labels werden übersprungen. So
        entstehen \num{347085} gelabelte Bilder von 262 Personen auf iPhones
        und iPads (\code{labels.csv}). Es gibt zwei Splits, jeweils 80/20:
        \code{norm\_subject} (\SI{80}{\percent} der Personen im Training,
        \SI{20}{\percent} im Test; keine Person in beiden) und
        \code{norm\_random} (zufällig nach Gruppen aus Person und Zielpunkt;
        dieselbe Person kann in Training und Test vorkommen). Für die Experimente
        werden \num{1000} Trainings- und \num{200} Testbilder zufällig gezogen,
        dazu 200 Validierungsbilder einer weiteren Person.
  \item \textbf{Eigene Aufnahmen} für die Personalisierung. Eine Person (der
        Autor) hat 27 Kalibriersitzungen mit zwei Laptop-Webcams aufgenommen.
        18 davon werden ausgewertet (\cref{sec:meth_aufnahme}).
\end{itemize}

\section{Aufnahme der Personalisierungsdaten}
\label{sec:meth_aufnahme}
Ein Aufnahmeprogramm zeigt einen roten Punkt auf dem Bildschirm, und die
Person folgt ihm mit den Augen. Gleichzeitig nimmt die Webcam Bilder auf. Zu
jedem Bild wird die Position des Punktes gespeichert, und zwar 6 Labels pro
Sekunde. Die Labels liegen in der Datei \code{norm\_labels.csv} jeder Sitzung;
die Positionen sind schon auf $[0,1]^2$ normiert. Es gibt drei verwendete
Generationen des Programms:
\begin{itemize}
  \item \textbf{Gen~1} (Sitzungen 01--04): Der Punkt bewegt sich gleichmäßig
        und ohne Sprünge über den ganzen Bildschirm.
  \item \textbf{Gen~2} (Sitzungen 05--09): Der Punkt springt alle 1,4 bis 2,7
        Sekunden an eine zufällige neue Stelle. Die Aufnahmen sind kurz (etwa
        \SI{30}{\second}).
  \item \textbf{Gen~4} (Sitzungen 19--27): Die Aufnahme besteht aus Abschnitten
        von \SI{5}{\second}. Jeder Abschnitt beginnt mit \SI{0,5}{\second}
        \code{WAIT} an einem neuen Startpunkt; danach bewegt sich der Punkt
        \SI{4,5}{\second} lang (\code{MOVE}). \SI{90}{\percent} der
        Startpunkte liegen in den Ecken. Das Programm speichert die Phase zu
        jedem Bild.
\end{itemize}
Gen~3 (Sitzungen 10--18) wurde verworfen, weil das Modell daraus nichts lernen
konnte. Die Sitzungen unterscheiden sich außerdem in Licht, Hintergrund,
Brille, Dauer (90 oder \SI{120}{\second} in Gen~4) und im Gerät. Es wurden
zwei Laptops mit jeweils eingebauter Webcam verwendet: \emph{Gerät~A}, ein
Microsoft Surface Pro~9 (Bildschirmbereich \num{1430}~$\times$~\num{880}
Pixel), für die Sitzungen 01--03 und 25--27, und \emph{Gerät~B}, ein
HP-Laptop (\num{1870}~$\times$~\num{1000} Pixel), für die Sitzungen 05--09
und 19--24. Für Sitzung 04 ist das Gerät nicht
protokolliert. Die Bilder werden auf beiden Geräten mit
\num{640}~$\times$~\num{480} Pixeln gespeichert.
\Cref{tab:session_conditions} zeigt alle Bedingungen.

\section{Vorverarbeitung}
\label{sec:meth_vorverarbeitung}
\begin{itemize}
  \item \textbf{Zielwerte:} Die Blickziele von GazeCapture werden durch Breite
        und Höhe des jeweiligen Bildschirms geteilt. So liegen alle Ziele im
        Einheitsquadrat $[0,1]^2$, unabhängig vom Gerät
        (\cref{sec:normalization}).
  \item \textbf{Bilder:} Das ganze Kamerabild wird auf
        $224\times224$ Pixel skaliert und mit den Mittelwerten und
        Standardabweichungen von ImageNet normiert. Es gibt keinen Gesichts-
        oder Augenausschnitt.
  \item \textbf{Bereinigte Labels:} Für Gen~4 kann man die \code{WAIT}-Bilder
        entfernen. Das Skript \code{session\_design.py --write-clean} schreibt
        dafür die Datei \code{norm\_labels\_clean.csv}.
\end{itemize}

\section{Modell und Training}
\label{sec:meth_modell}
Das Modell besteht aus zwei Teilen:
\begin{enumerate}
  \item einem \textbf{ResNet-18}~\cite{he2016resnet}, auf ImageNet
        vortrainiert (\code{ResNet18\_Weights.DEFAULT}), als Feature
        Extractor. Es liefert 512 Features pro Bild.
  \item einem \textbf{Regression Head}: ein MLP aus vier vollständig
        verbundenen Schichten $512 \to 512 \to 256 \to 128 \to 2$ mit ReLU,
        Dropout $p=0,2$ nach der zweiten Schicht und einer Activation am
        Ausgang (Standard: Sigmoid). Er berechnet daraus die Position $(x,y)$.
\end{enumerate}
Trainiert wird mit L1-Loss (\code{nn.L1Loss}), AdamW (Learning Rate $10^{-4}$,
Weight Decay $10^{-5}$), Batch Size 32 und Early Stopping (Abbruch nach drei
Epochs ohne Verbesserung, höchstens 500 Epochs). Das Vorgehen hat zwei Stufen:
\begin{enumerate}
  \item \textbf{Basismodell (E0):} Der Head wird auf GazeCapture trainiert
        (\code{train\_base.py}). Von fünf Seeds wird der Lauf mit dem kleinsten
        Validierungsfehler behalten.
  \item \textbf{Personalisierung (E1--E3):} Der Head des Basismodells wird mit
        den Bildern einer eigenen Sitzung weiter trainiert
        (\code{personal\_exp.py}). Der Backbone bleibt dabei eingefroren,
        auch die BatchNorm-Statistiken.
\end{enumerate}
Weil der Backbone eingefroren ist, werden seine Features nur einmal berechnet
und gespeichert (Feature-Cache). Danach wird nur der Head trainiert. Ein Lauf
dauert deshalb nur Sekunden.

\section{Auswertung}
\label{sec:meth_evaluation}
\begin{itemize}
  \item \textbf{Fehlermaß:} Diagonalfehler $\Ediag$, also der euklidische
        RMSE in Prozent der Bildschirmdiagonale (\cref{eq:ediag}).
  \item \textbf{Baselines:} Die Konstant-Baseline sagt immer den Mittelwert der
        Trainingsziele vorher, die Zufalls-Baseline rät zufällig. Beide werden
        auf denselben Testdaten wie das Modell berechnet
        (\cref{sec:baselines}).
  \item \textbf{Statistik:} Bei GazeCapture gibt es fünf Wiederholungen pro
        Einstellung mit \SI{98}{\percent}-Confidence Interval (Student-$t$)
        und dem Mann-Whitney-$U$-Test. Bei den Sitzungen nutze ich den
        Wilcoxon-Vorzeichen-Rang-Test (gepaart) und den Kruskal-Wallis-Test
        (Vergleich von Gruppen).
  \item \textbf{Aufteilung der Sitzungsdaten:} zufällig 70/15/15, sequenziell
        (zeitliche Blöcke mit einer Lücke von einer Sekunde) und zwei Videos
        (Training auf einer Sitzung, Test auf einer anderen)
        (\cref{sec:loading}).
  \item \textbf{Räumliche Analyse:} Fehler in Mitte, Rändern und Ecken des
        Bildschirms (\cref{sec:partition}).
\end{itemize}

\section{Implementierung}
\label{sec:meth_implementierung}
Der Code ist in Python mit PyTorch und torchvision geschrieben. Die wichtigsten
Teile sind:
\begin{itemize}[nosep]
  \item \code{main.py}: lädt GazeCapture, normiert die Ziele und erstellt die
        Splits;
  \item \code{model.py}, \code{dataset.py}: Modell, Losses, Training,
        Laden der Daten und Feature-Cache;
  \item \code{metrics.py}: Diagonalfehler und Baselines (eine einzige Quelle
        für alle Skripte);
  \item \code{exp.py}: Hyperparameterstudie mit fünf Wiederholungen;
  \item \code{train\_base.py}, \code{personal\_exp.py}: Basismodelle und
        Personalisierung (E0--E3);
  \item \code{session\_design.py}: Abdeckung, Transition Frames und bereinigte
        Labels;
  \item \code{spatial\_analysis.py}, \code{plot\_targets.py}: Abbildungen;
  \item \code{auswertung.py}: erstellt die Tabellen aus den Ergebnisdateien;
  \item Notebook \code{notebooks/Masterarbeit.ipynb}: führt alle Experimente aus und
        erstellt Tabellen und Abbildungen.
\end{itemize}
Alle Experimente laufen ohne Grafikkarte auf der CPU des HP-Laptops
(Gerät~B): AMD Ryzen~7 3700U, etwa \SI{13}{\giga\byte} nutzbarer
Arbeitsspeicher, Ubuntu~26.04.1~LTS, Python~3.12.3, PyTorch~2.12.1 und torchvision~0.27.1
(jeweils CPU-Version). Die Befehle zum Wiederholen der
Experimente stehen im Anhang (\cref{sec:anhang_befehle}). Der Code und die
Web-Demo sind auf GitHub verfügbar:
\url{https://github.com/RafeiKAr/gaze-estimation}. Eine Kopie liegt in der
BHT-Cloud: \url{https://cloud.bht-berlin.de/index.php/s/P4GkQ9J8zWHbdf9}. Die
Bilder und Videos der eigenen Aufnahmen zeigen das Gesicht des Autors. Sie sind
deshalb nicht öffentlich, sondern liegen in einem privaten Dataset auf
Hugging Face.


